% Izvor: PDF strana 6, gornji slajd. Redizajn B / etf-v1.
\slajdid{01-s011}
\begin{frame}[label=01-s011]{Binarna promenljiva}
% element: 01-s011:e01
Umesto „netačno“ i „tačno“ koristimo \textbf{logičku nulu} \(0\) i \textbf{logičku jedinicu} \(1\).
\[\boxed{A=\{0,1\},\qquad P\in A}\]
% element: 01-s011:e02
\vspace{4mm}
\Rezultat{Pridružujemo promenljivu \(P\) stanju prekidača:\par\vspace{3mm}
\begin{minipage}{.48\linewidth}\textbf{Zatvoren prekidač}\quad\(P=1\)\end{minipage}
\hfill
\begin{minipage}{.48\linewidth}\textbf{Otvoren prekidač}\quad\(P=0\)\end{minipage}}
% element: 01-s011:e03
\par\vspace{6mm}
\Rezultat{\(P\) je binarna promenljiva: može imati samo dva stanja, \(0\) ili \(1\).}

\end{frame}
